\sgaStableTarget{sga6.e06.lem.1-10}\textbf{Lemma 1.10.} \uline{With the notation and hypotheses of} 1.1, \uline{let} @@M0@@ and @@M1@@ \uline{be two locally free sheaves of finite type on} @@M2@@, \uline{such that}:
@@M3@@
\uline{then there exists an integer} @@M4@@, \uline{such that for} @@M5@@, \uline{we have}:
@@M6@@

Let @@M7@@ be the free abelian group generated by the classes (up to isomorphism) of locally free sheaves of finite type on @@M8@@. To say that, in @@M9@@, @@M10@@ and @@M11@@ have the same class is to say, by the construction of @@M12@@, that there exist exact sequences of @@M13@@-Modules locally free of finite type:
@@M14@@
and integers @@M15@@, such that in @@M16@@ we have the equality:
@@M17@@
For every @@M18@@, tensoring with @@M19@@ defines an automorphism of @@M20@@, whence:
@@M21@@
the sequences @@M22@@ remaining exact.

Let @@M23@@ be the free abelian group generated by the classes (up to isomorphism) of quasi-coherent sheaves on @@M24@@, and @@M25@@ the subgroup generated by the classes of locally free sheaves of finite type on @@M26@@. Direct image defines a homomorphism @@M27@@, and we therefore still have:
